% Clase y configuracion de tipo de documento
\documentclass[10pt,a4paper,spanish]{article}
% Inclusion de paquetes
\usepackage{a4wide}
\usepackage{amsmath, amscd, amssymb, amsthm, latexsym}
\usepackage[spanish]{babel}
\usepackage[utf8]{inputenc}
\usepackage[width=15.5cm, left=3cm, top=2.5cm, height= 24.5cm]{geometry}
\usepackage{fancyhdr}
\pagestyle{fancyplain}
\usepackage{listings}
\usepackage{enumerate}
\usepackage{xspace}
\usepackage{longtable}
\usepackage{caratula}
\usepackage{hyperref}
\usepackage{graphicx}
\graphicspath{{img/}}

% Encabezado
\lhead{Organización del Computador II}
\rhead{Grupo: Fuga Villera Nro. 2}
% Pie de pagina
\renewcommand{\footrulewidth}{0.4pt}
\lfoot{Facultad de Ciencias Exactas y Naturales}
\rfoot{Universidad de Buenos Aires}

\begin{document}

% Datos de caratula
\materia{Organización del Computador II}
\titulo{Trabajo Práctico 2: Experimentos}
%\subtitulo{}
\grupo{
	Grupo: Fuga Villera Nro. 2
	\href{https://www.youtube.com/watch?v=wBNjDRXJNyY}{
		\includegraphics[height=0.4cm,keepaspectratio]{YouTube-icon-full_color.png}
	}
}

\integrante{Gabriel Matles}{397/12}{gabriel29m@gmail.com}
\integrante{Manuel Mena}{313/14}{manuelmena1993@gmail.com}
\integrante{Francisco Demartino}{348/14}{demartino.francisco@gmail.com}

\maketitle

\newpage

% Para crear un indice
\tableofcontents

% Forzar salto de pagina
\clearpage

\section{Objetivos generales}

El objetivo de este trabajo práctico es evaluar la eficiencia del modelo de programación SIMD mediante la implementación de diversos algoritmos en lenguaje Ensamblador utilizando instrucciones SSE. \\
\indent Las mediciones se realizan mediante pruebas empíricas del código frente a algoritmos que cumplen la misma especificación, implementados en un lenguaje de alto nivel (C). \\
\indent En este proyecto, los algoritmos a implementar se basaron en el procesamiento de imágenes, en el cual el uso del modelo SIMD es provechoso.

\section{Calidad de las Mediciones}
Antes de medir tiempos y sacar conclusiones de los mismos, es importante tener un esquema básico de como se realizarán las mediciones en el presente trabajo, y bajo que condiciones. \\

\textbf{Características de la computadora en la que los experimentos fueron realizados:} \\
\textit{Procesador:} Intel Core iX XGH X CPUs @ XGHz\\
\textit{Memoria RAM:} X GB \\
\textit{Sistema operativo:} Ubuntu 14.04 64 bits. \\

Las mediciones se realizaron corriendo los filtros X veces cada uno. Se realizó un promedio, y estos promedios son los que se pueden apreciar en los distintos gráficos. Los tiempos están medidos en ¿¿¿¿¿¿¿ciclos de clock??????.

\section{Experimentación}



\section{Observaciones}
	\includegraphics[keepaspectratio]{gen/diff.png}
	\begin{enumerate}
		\item un item
		\item otro item
	\end{enumerate}

% Otro salto de pagina
% \newpage

\section{Resolución}


% \subsection{Ejercicio X}

\subsection{Auxiliares}

\end{document}
